\section{Pass a search request without hiding its state}
Implement \texttt{h263\_search16}. Its first argument, \texttt{\$a0}, points to
an eight-byte-aligned request containing eight native-endian, unsigned 64-bit
fields. The second, \texttt{\$a1}, points to an eight-byte-aligned, writable
24-byte result area. Native endian means the simulator's selected target byte
order; these in-memory records are not H.263 file syntax.
\begin{center}
\tagpdfsetup{table/header-rows={1}}
\begin{tabular}{rll}
\textbf{Byte offset} & \textbf{Field} & \textbf{Meaning or range}\\
0 & Current pointer & Top-left sample of the current block\\
8 & Reference pointer & Top-left sample of the reference plane\\
16 & Current stride & 16 through 4096 bytes\\
24 & Reference stride & 176 through 4096 bytes\\
32 & $x$ & Current picture coordinate, 0 through 160\\
40 & $y$ & Current picture coordinate, 0 through 128\\
48 & Radius & 0, 1 or 2 full pixels\\
56 & Comparison mode & 0 scalar; 1 MSA
\end{tabular}
\end{center}
The caller supplies mapped records, mapped current samples and all 176 active
samples in every row of the reference plane. The result area must not overlap
the request or either input. The two read-only inputs may overlap each other.
They may have different strides, even if their starting pointers are equal.

Validate record alignment and non-wrapping record extents before reading the
request. Validate all fields, the current minimal extent $15s_C+16$ and the
reference minimal extent $143s_R+176$ before reading any pixels. The larger
reference extent describes the plane supplied by this interface; it does not
authorize reading its padding. An exclusive end that wraps is invalid. Stride,
coordinate, radius and mode checks must examine the whole 64-bit field, so a
large value cannot pass because its low word resembles a small valid value.

On success, return \texttt{\$v0}=0 and the number of unique admitted candidates
in \texttt{\$v1}. Write signed 64-bit $d_x$ at result offset zero, signed
64-bit $d_y$ at offset eight, and unsigned 64-bit SAD at offset sixteen.
On invalid arguments, return one and zero without reading pixels or changing
any result byte. Request reads needed for validation are allowed; reading a
pixel speculatively before validation is not. Mapped-memory preconditions are
supplied by the caller, rather than discovered by probing arbitrary addresses.

Preserve saved general registers, the stack pointer, the request and both
inputs. Scratch registers and vector registers may change. A nested SAD call
can overwrite the return-address register and its argument registers. Store
long-lived search state in saved registers or a valid stack frame, and restore
the caller's state before returning. The public harness provides 16\,384 bytes
of private stack; it does not prescribe your frame size or local labels.

The ordinary Run/Debug caller is
\path{drafts/h263-motion-construction/main.asm}. Carry your completed SAD entries
into this new working source; retain one ordinary main entry. The supplied search
and SAD bodies return development status four until replaced. The caller installs
its two pointer fields at runtime, then calls the search. Its named data fields
retain the returned status and candidate count for inspection.

Prepare a worked copy from your current source, preserving its instructions:
\begin{verbatim}
python3 drafts/h263-motion-construction/prepare.py \
    --source my-search.asm --mode 1 \
    --output worked-search.asm
\end{verbatim}
Mode zero chooses scalar comparisons. The completed worked caller prints
\texttt{-2 -2 2976} and exits successfully. Before completion, its development
message and failing exit status are expected. At the breakpoint
\texttt{motion\_caller\_return}, before the first store sequence executes, the
result record is already complete and the returned registers are zero and 25.
The caller has not yet copied those registers to its named status/count fields.
